\scanpage{8}

% loc: top, continuation
tel qu'on ait un isom fonctoriel en $X$

% loc: first Hom display
\[\operatorname{Hom}_{E}(X,\xi)\simeq\operatorname{Hom}_{E_1}(\alpha'X,\xi)\quad(\simeq\operatorname{Hom}_{E}(X,\alpha\xi))\]

% loc: question
en d'autres termes a-t-on $\xi\in\operatorname{Ob}\overline{E_1}$ (image essentielle de $\alpha$) ? Or quelle est la relation entre $\Sigma$ et $E_1$ ?

% loc: equivalences
\[
\begin{aligned}
\deleted{\unsure{u\in\Sigma}}\ (u:X\to Y)\in\Sigma
&\Longleftrightarrow\alpha'(u)\inserted{:\alpha'(X)\simeq\alpha'(Y)}\ \text{un isom}
\Longleftrightarrow\underset{\text{pour }\forall\, Z_1\in\operatorname{Ob}E_1}{\operatorname{Hom}_{E_1}(\alpha'(Y),Z_1)\simeq\operatorname{Hom}_{E_1}(\alpha'(X),Z_1)}\\
&\Longleftrightarrow\operatorname{Hom}_{E_1}(Y,\alpha Z_1)\simeq\operatorname{Hom}_{E_1}(X,\alpha Z_1)\quad\text{pour tt } Z_1\in\operatorname{Ob}E_1,
\end{aligned}
\]

% loc: last paragraph
Donc la question est si on récupère bien $\overline{E_1}$ comme sous-catégorie str. pleine de $E$ comme la sous catégorie \inserted{str.} \deleted{\unsure{pleine}} $\Phi(\Sigma)\subset E$ formée des $\xi$ \unsure{tel|tels} que $\operatorname{Hom}(Y,\xi)\xrightarrow{\sim}\operatorname{Hom}(X,\xi)$ pour tt $u:X\to Y$ dans $\Sigma$. Or c'est vrai sauf erreur \ldots.

% archive sheet number: 7
